\section{Review Protocol}
\label{sec:review-protocol}

\subsection{Review Questions and Scope}

This survey asks: (i) how are robot courses represented and generated; (ii)
which constraints, validity checks, and evaluation measures make generated courses
usable for learning and control; and (iii) how can these choices be reported so
that results can be compared and reproduced? The unit of technical interest is
an output spatial layout that constrains traversal: a track or corridor, road,
ordered gate or waypoint sequence, buoy course, or world layout. We consider
such outputs for ground, aerial, and maritime robots, including simulated systems
when their vehicle model and control task make course geometry consequential.
Transferable adjacent methods are eligible when they contribute a representation,
generator, constraint, metric, feasibility test, interface, or benchmark that
maps to this unit without changing the contribution's essential claim.

The scope is deliberately broader than the historical state-of-the-art document,
which focused on track- and road-like geometry generation. It also covers
course distributions for robot learning and control, gate and waypoint courses,
simulation and benchmark interfaces, official specifications and competition
materials, and adjacent procedural road, map, network, level, and test-input
generation. Conversely, a racing-line planner, controller, traffic generator,
or appearance/dynamics randomizer is not a course-generation method merely
because it operates on a track.

\subsection{Discovery, Reports, and Works}

Discovery combines a documented historical bootstrap with independent query
streams for ground systems, aerial and maritime systems, geometry and
reinforcement-learning citation expansion, and simulator, benchmark, and
official-artifact interfaces. Queries span direct robot-racing terminology as
well as procedural road generation, autonomous-vehicle testing, racing-game
procedural content generation, and portable course formats. Discovery is used
to assemble a candidate population, not as evidence for a technical claim;
primary reports and authoritative artifacts are assessed during screening.

The screening unit is a \emph{report}: one identifiable article, preprint,
thesis, technical report, standard, dataset, release, repository revision, or
official document. Exact duplicate reports are collapsed with their aliases
retained. Versioned or companion reports remain linked until their unique
evidence has been assessed. The synthesis unit is the underlying \emph{work}:
linked reports are grouped by stable identifiers, authorship, artifact identity,
and contribution overlap; a canonical citable report is selected, and a work is
counted at most once in any later quantitative synthesis.

\subsection{Two-Pass Evidence Procedure}

Pass~1 is a binary, evidence-based retention screen. A report is retained only
when the frozen packet establishes at least one of: a direct course-generation
contribution; a fixed-course property that transfers to a survey or benchmark
design decision; or contextual evidence from an inspected survey or systematic
review. Excluded reports receive one controlled, source-specific exclusion
criterion. Pass~2 is separate multi-label contribution coding, which assigns
evidence tier and supported method, representation, interface, dataset,
benchmark, validity-test, and metric labels. Thus a retained-report count is
neither a method count nor a prevalence estimate.

Fixed routes, static competition layouts, and contextual works are retained only
when they provide citable transferable representation, interface, benchmark,
evaluation, or terminology context. They are not relabelled as generators.
In particular, running, optimizing, or evaluating control on a fixed route does
not establish a course-generation contribution.

Material evidence is ranked as follows: inspected primary reports and
authoritative companion artifacts support technical claims; version-identified
official documentation may support artifact, interface, or benchmark claims;
and directly inspected reviews may support only field, terminology, or
literature-gap context. Each deciding item requires a canonical source URL,
precise page, section, figure, algorithm, repository-path/line, or stable
documentation anchor, plus the inspected version, retrieval date, and frozen
artifact or archive locator. The packet also records the associated artifact
identity where bytes are available. Titles, search snippets, and unverified
abstract inferences are not material evidence; an abstract-only record cannot
be included.

\subsection{Independent Screening and Adjudication}

Each report is rated independently twice against the same versioned,
integrity-bound evidence packet. The calibration sample was screened before
main screening under the final controlled status and criterion vocabulary; the
main phase was released only after the calibration gate. Reviewers do not see
the paired rating before both ratings are locked. Adjudication is append-only
and is triggered cumulatively by A1, a status disagreement; A2, a criterion
disagreement; A3, unequal normalized exclusion-reason text despite otherwise
matching exclusions; or A4, a pre-existing unresolved conflict in the frozen
record. These triggers distinguish disagreements about eligibility from
differences in required explanatory text.

\begin{table}[t]
  \centering
  \caption{Audited Pass~1 screening reliability and adjudication workload. These
  are process statistics, not final corpus results.}
  \label{tab:screening-checkpoint}
  \small
  \begin{tabularx}{\linewidth}{@{}lY@{}}
    \toprule
    Measure & Audited value \\
    \midrule
    Candidate reports & 202 \\
    Calibration status agreement & 27/30 (90.0\%) \\
    Main status agreement & 170/172 (98.84\%) \\
    Pooled status agreement & 197/202 (97.52\%) \\
    Pooled exact status--criterion agreement & 195/202 (96.53\%) \\
    Cumulative adjudication triggers & 132; 124 were normalized reason-text
    mismatches, rather than eligibility disagreements \\
    \bottomrule
  \end{tabularx}
\end{table}

\subsection{Current Non-Final Checkpoint}

\Cref{tab:screening-checkpoint} summarizes the audited screening process for
202 reports. The current adjudication workbook is explicitly non-final: 75
reports are draft-included, 125 are draft-excluded, and two are held. Among
the draft inclusions, 53 use \texttt{full\_text} evidence and 22 use
\texttt{official\_documentation}; none uses \texttt{abstract\_only} evidence.
The 132 cumulative triggers should not be read as 132 eligibility disagreements:
most arose from the protocol's intentionally sensitive A3 comparison of
normalized reason text.

Two reports, \texttt{C0046} and \texttt{C0173}, remain held because their frozen
PDFs were structurally unreadable. They require either a corrected, re-frozen
packet followed by rerating or a governed corrupt-artifact disposition. They
are therefore not silently treated as ordinary exclusions or inclusions.

\subsection{Audit Trail and Limitations}

The review preserves the query matrix and search ledger, report and evidence
manifests, frozen reviewer releases, locked results, calibration decision,
adjudication workbook, and versioned evidence locators. This audit trail makes
the evidence basis and each later transformation inspectable, but it does not
remove limitations: discovery may miss eligible work, access limitations can
prevent assessment, versioned online artifacts can change, and automated
structural checks cannot replace accountable scholarly verification. No final
corpus, prevalence, or performance claim will be made until Pass~2 coding,
citation-key projection, execution receipts, and accountable-author
verification are complete.
